\section{Conclusions}
\label{sec:conclusions}

\red{This paper presented an AI-driven multi-source urban analytics framework
for people forecasting in the historic borough of Dozza. The framework combines
pedestrian camera counts, mobile-network indicators, weather observations, and
event-related information in a single hourly modeling pipeline. The updated
evaluation uses causal temporal baselines and an embargo between train and test,
so that the reported forecast errors do not rely on future target information.}

\red{The strongest empirical result concerns pedestrian-flow forecasting.
Feature-based tabular models, especially Extra Trees, consistently outperform
the strict last-hour baseline for entrances and exits across the evaluated
horizons. This indicates that recent pedestrian dynamics and contextual
variables provide useful information for visitor-flow prediction in a small
historic destination.}

\red{The TIM-derived targets show a more nuanced pattern. On the final
chronological test split, feature-based models often improve over last-hour
persistence for nationality and age-band targets. However, one-hour rolling
validation still selects last-hour persistence for those TIM targets. The
appropriate conclusion is therefore not that complex models always dominate,
but that their advantage depends on the target, horizon, and validation
protocol.}

\red{The interpretability analyses support this view. Pedestrian-history
features are most important for pedestrian-flow models. TIM and target-history
features dominate the demographic models. Event variables are selected and can
rank highly in SHAP summaries for some targets, but their ablation effect is
mixed. They should be treated as exploratory signals rather than as stable
drivers of forecast accuracy.}

\red{The study has three main limitations. First, the analysis is limited to
the period where all data sources overlap. Second, it studies one borough, so
external validity must be tested on other small tourism destinations. Third,
the event dataset is still incomplete and does not provide reliable measures of
attendance, media impact, or visitor attractiveness. Future work should extend
the observation period, improve the event representation, and replicate the
framework in other historic towns before operational deployment in real-time
decision-support systems.}
